\section{Arquitetura Distribuída e Protocolos de Comunicação}
\label{sec:arquitetura_distribuida}

O sistema proposto foi modelado sob um paradigma cliente-servidor fracamente acoplado, dividindo a carga de trabalho entre processamento de borda (\textit{edge computing}) e armazenamento especializado em nuvem. A separação de responsabilidades visa otimizar o uso de recursos locais e garantir a escalabilidade horizontal do banco de dados de vetores.

\subsection{Topologia da Rede e Desacoplamento de Nós}
A topologia do sistema é composta por dois nós principais que operam de forma geograficamente distribuída:

\begin{itemize}
    \item \textbf{Nó Cliente (Local):} Executa em ambiente computacional próprio, encarregado das tarefas de computação intensiva na origem, que compreendem a leitura do sistema de arquivos local, a decodificação de matrizes de imagens e a geração dos vetores de características de 512 dimensões.
    \item \textbf{Nó Servidor (Nuvem):} Consiste em um cluster gerenciado do \textit{Qdrant Cloud}, hospedado na infraestrutura da \textit{Amazon Web Services} (AWS) na região da América do Sul (\textit{sa-east-1}). Este nó é responsável exclusivo pela indexação espacial, persistência e processamento de consultas de similaridade de alta dimensionalidade.
\end{itemize}

A comunicação entre as extremidades ocorre por meio de requisições HTTP utilizando uma interface de programação de aplicações baseada na arquitetura REST (\textit{RESTful API}), implementada sobre as bibliotecas de transporte \texttt{httpx} e \texttt{httpcore}.

\subsection{Base de Dados de Validação e Protocolo de Experimentos}
A validação experimental da infraestrutura distribuída foi conduzida utilizando a base de dados de faces da FEI (\textit{FEI Face Database})~\cite{thomaz2006fei}. Este repositório biométrico é composto por 2.800 imagens coloridas com resolução original de $640 \times 480$ pixels, obtidas a partir de um grupo de 200 indivíduos, apresentando uma divisão equitativa de 100 sujeitos do sexo masculino e 100 do sexo feminino, com idades variando entre 19 e 40 anos. Cada participante possui 14 registros fotográficos capturados contra um fundo branco homogêneo, abrangendo diferentes expressões faciais e rotações de perfil de até cerca de 180 graus, com variações de escala estimadas em 10\%.

No modelo de testes aplicado ao sistema distribuído, o conjunto de dados foi segmentado e alocado logicamente para simular um cenário real de identificação remota:

\begin{itemize}
    \item \textbf{Fase de Ingestão:} O subconjunto composto pelas imagens em posição estritamente frontal e com expressão neutra permaneceu no nó cliente para o processo de extração local e alimentação síncrona inicial da coleção hospedada no \textit{Qdrant Cloud}.
    \item \textbf{Fase de Consulta:} As imagens que apresentavam variações de pose, rotação e expressões do mesmo indivíduo foram isoladas em um diretório de testes secundário no nó local. Esses arquivos foram utilizados para alimentar o laço contínuo de requisições externas, validando a capacidade do método remoto \texttt{query\_points} em correlacionar vetores com distorções geométricas ao vetor original armazenado na nuvem.
\end{itemize}

\subsection{Implementação com a leitura de Rostos em Tempo Real}
